\section{Introduction}
\label{sec:introduction}

\subsection{The result}

A \emph{bending} of a closed surface is a continuous family of
isometric embeddings of one fixed Riemannian metric into $\R^3$. It is
nontrivial when the shape of the surface changes, that is, when the
images are not all congruent. Connelly's flexible sphere
\cite{Connelly1977} shows that nontrivial bendings exist for embedded
polyhedra. Problem~1.1 of Ghomi's list \cite{Ghomi2019}, which is
attributed there to Euler and Maxwell, asks whether a closed $C^2$
surface in $\R^3$ can be flexible. We construct a $C^\infty$ example.

\begin{maintheorem}\label{thm:main}
There are a smooth surface $\Sigma$ diffeomorphic to $S^2$, a smooth
Riemannian metric $g_\Sigma$ on it, and a jointly smooth map
$\mathcal F:\Sigma\times\R\to\R^3$ such that the maps
$F_\tau=\mathcal F(\,\cdot\,,\tau)$ satisfy
\[
 F_\tau\text{ is an embedding},\qquad
 F_\tau^*g_{\mathrm E}=g_\Sigma,\qquad F_{\tau+1}=F_\tau
 \qquad\text{for all }\tau\in\R,
\]
where $g_{\mathrm E}$ is the Euclidean metric. At least two of the
images $F_\tau(\Sigma)$ are not congruent under the Euclidean isometry
group. An inner disk and a complementary cap of $\Sigma$ are fixed
pointwise throughout the family. The area, the enclosed volume and the
total oriented mean curvature of $F_\tau(\Sigma)$ do not depend on
$\tau$.
\end{maintheorem}

Three conventions are implicit in this statement. A \emph{closed}
surface is compact and without boundary. \emph{Nontrivial} always
refers to noncongruence of the images: relabelling the points of a
fixed surface by a family of intrinsic isometries moves material
points but does not change the shape, and is not a bending in our
sense. The mean curvature is $(k_1+k_2)/2$ for the outward normal.

\subsection{The mechanism}
\label{subsec:mechanism}

The family in Theorem~\ref{thm:main} is completely explicit once one
particular disk is available. Let $D$ be a disk with a rotationally
symmetric metric $g$, and write $\varrho_\vartheta$ for its intrinsic
rotation by the angle $\vartheta$. Suppose $F:D\to\R^3$ is an isometric
embedding. If $R_\vartheta$ denotes the rotation of $\R^3$ about the
vertical axis, then
\begin{equation}\label{eq:phase-family}
 F_\tau=R_{-2\pi\tau}\circ F\circ\varrho_{2\pi\tau}
\end{equation}
is again an isometric embedding of $(D,g)$ for every $\tau$, and its
image is the rotated copy $R_{-2\pi\tau}F(D)$ of the original image.
Two cases are uninteresting. If $F(D)$ is a surface of revolution,
the image does not move at all. If $F(D)$ is not a surface of
revolution, the image does move, but in general so does its boundary
curve, and nothing can be attached to it.

There is a third possibility. Suppose that $F(D)$ is not a surface of
revolution, but that near the boundary of $D$ the embedding is the
inclusion of a flat horizontal annulus, $F(x)=(x,c)$. There the two
rotations in \eqref{eq:phase-family} cancel, so $F_\tau=F$ near the
boundary for every $\tau$. All the maps $F_\tau$ can then be glued to
one and the same complementary cap. The result is a closed surface in
which a wavy annulus turns relative to a stationary shell
(Figure~\ref{fig:surface-motion}). If the shell carries a mark which
breaks its rotational symmetry, no rigid motion can account for the
relative rotation, and the shape of the closed surface changes.

\begin{figure}[htbp]
\centering
\begin{tikzpicture}[scale=.91,>=Stealth,font=\small]
\foreach \shift/\phase/\word in {0/0/initial,6.6/55/later}{
 \begin{scope}[xshift=\shift cm]
 \fill[gray!15] (0,0) circle (2.05);
 \fill[blue!6] (0,0) circle (1.65);
 \draw[gray!65,thick] (0,0) circle (2.05);
 \draw[dashed] (0,0) circle (1.65);
 \foreach \radius in {.85,1.12,1.38}{
  \draw[blue!70!black,thick,rotate=\phase]
   plot[domain=0:360,samples=65,variable=\t]
   ({(\radius+.10*cos(2*\t))*cos(\t)},
    {(\radius+.10*cos(2*\t))*sin(\t)});}
 \filldraw[fill=green!15,draw=green!45!black] (0,0) circle (.52);
 \fill[red!75!black] (1.30,1.33) circle (2.4pt);
 \node at (0,-2.36) {\word\ position};
 \end{scope}}
\draw[->,thick] (2.65,0)--(3.95,0);
\node[above] at (3.3,0) {phase $\tau$};
\node[below,align=center] at (3.3,-.1) {same\\metric};
\node[above] at (0,2.15) {fixed inner disk};
\draw[->,green!45!black] (-.5,2.0)--(-.18,.45);
\node[above] at (6.6,2.15) {fixed marker on the cap};
\draw[->,red!75!black] (7.4,2.0)--(7.90,1.36);
\end{tikzpicture}
\caption{The relative rotation, seen from above. The blue curves are
schematic height contours of the moving annulus; they are not numerical
data for the constructed surface. The dashed circle bounds the planar
collar, which stays fixed pointwise, as does the green inner disk. The
red marker on the cap also stays fixed while the nonradial part turns,
so no ambient rotation accounts for the whole motion.}
\label{fig:surface-motion}
\end{figure}

The motion itself therefore requires no analysis; in particular, no
infinitesimal bending has to be integrated in the variable $\tau$.
Everything rests on the existence of the disk, and the difficulty lies
in the requirement that it be \emph{exactly} planar on a collar of its
boundary while being nonradial inside.

\subsection{The analytic core}

Such a disk is produced as follows. Fill the middle of the disk by a
flat piece, so that only an annulus $(r_-,r_+)\times\T$ remains, where
$\T=\R/2\pi\Z$. On the annulus take the metric of a radial graph with
slope $w$,
\begin{equation}\label{eq:intro-metric-form}
 g=(1+w(r)^2)\,dr^2+r^2\,d\theta^2,
\end{equation}
and look for an isometric embedding of the form $(P,z)$, where $z$ is
a height function and $P$ takes values in the horizontal plane. Such
an embedding exists locally exactly when the \emph{remainder metric}
$g-dz\otimes dz$ is positive definite and flat, and flatness is a
single scalar equation for $z$, the Darboux equation. The disk of
Section~\ref{subsec:mechanism} is then obtained from the following
statement, which is the main analytic result of the paper.

\begin{maintheorem}\label{thm:annulus-intro}
There exist radii $0<r_-<r_+$, a smooth function $w>0$ on $(r_-,r_+)$
and a smooth function $z$ on $(r_-,r_+)\times\T$ with the following
properties, where $g$ is the metric \eqref{eq:intro-metric-form}.
\begin{enumerate}
\item $|dz|_g<1$, and the metric $g-dz\otimes dz$ is flat.
\item $z(r,\theta+\pi)=z(r,\theta)$, and $z$ is not independent of
 $\theta$.
\item $w$ extends smoothly by zero across $r_-$ and $r_+$, and $z$
 extends smoothly by a constant across each of them.
\end{enumerate}
\end{maintheorem}

Condition~(iii) says that $w$, and $z$ minus its boundary value, vanish
to infinite order at both boundary circles. After a slight enlargement
of the annulus this gives open collars on which $w=0$ and $z$ is
constant, which is the exact planarity needed for the mechanism. A
small error at the boundary would be useless; this is why flatness to
infinite order is the right target. The precise form of
Theorem~\ref{thm:annulus-intro}, with quantitative information on $w$
and $z$, is Theorem~\ref{thm:new-annular-height}.

The passage from Theorem~\ref{thm:annulus-intro} to
Theorem~\ref{thm:main} is carried out in
Section~\ref{sec:geometric-closure} and uses only classical
differential geometry. The rest of the paper proves
Theorem~\ref{thm:annulus-intro}.

\subsection{Three difficulties, and how they are met}
\label{subsec:difficulties}

We look for the height in the form $z=h+qu$, where $h'=w$ is the radial
height whose graph has the metric $g$, and $q$ is a small parameter. To
first order in $q$ the Darboux equation is
\begin{equation}\label{eq:intro-linear}
 \mathscr L_wu=(wu_r)_r+\frac{w'}{r}\,u_{\theta\theta}=0 .
\end{equation}
A nonzero solution of this equation which is constant to infinite order
at both ends is an infinitesimal version of what we want, and the plan
is to choose $w$ so that one exists and then to correct it to a
solution of the nonlinear equation. Three features of the problem shape
everything that follows.

\paragraph{Flat ends force infinitely many changes of type.}
For $w>0$ the operator \eqref{eq:intro-linear} is elliptic where $w'>0$
and hyperbolic where $w'<0$. At a flat end $w$ tends to zero. If $w$
were monotone near that end, no nonzero angular mode of
\eqref{eq:intro-linear} could vanish there together with its flux
(Proposition~\ref{prop:obs-monotone-tail}). So $w$ must oscillate: near
each end there are infinitely many circles on which the equation
changes type, and they accumulate at the boundary. We build such a
profile from infinitely many copies of one explicitly solvable cell,
with lengths and amplitudes decaying fast enough for smoothness
(Sections~\ref{sec:rosen-reference-packet} and~\ref{sec:radial-tails}).

\paragraph{Two unrelated resonance conditions.}
On each hyperbolic annulus, characteristics turn by a fixed
\emph{travel angle} $T$, and solving along them involves dividing by
$\sin(jT)$ at the angular frequency $j$. Independently, the solutions
on neighbouring annuli must be matched, and the matching fails at
eigenvalue-like resonances. A finite model in which both sets of
resonances can be computed exactly (Section~\ref{sec:worked-example})
shows that neither condition implies the other. In the actual profile
the travel angles are quadratic irrationals fixed by the choice of the
cells, while matching resonances are avoided by adjusting the relative
heights of consecutive cells. The outcome is a two-sided inverse of
$\mathscr L_w$ on the whole annulus, in spaces weighted towards the
flat ends (Section~\ref{sec:radial-reference-inverse}).

\paragraph{The inverse has to survive perturbation.}
The Nash--Moser method needs an inverse of the linearization at every
state near the reference one, not only at the reference. For a nearby
state both the type-change circles and the characteristic return maps
move, so its linearization does not even act on the same space. We
remove this problem by subtracting from the linearization two terms,
each proportional to an equation that is not yet solved: one fixes the
type-change circles, the other fixes the return maps. Both terms vanish
at a solution, and they produce exactly the quadratic error allowed by
Hamilton's form of the implicit function theorem \cite{Hamilton1982}
(Sections~\ref{sec:fixed-geometry}--\ref{sec:solving-the-equation}).

\subsection{Organization of the proof}
\label{subsec:organization}

Figure~\ref{fig:proof-map} shows the logical structure, and
Table~\ref{tab:sections} lists what each section contributes. The
paper has three parts.

\begin{figure}[htbp]
\centering
\begin{tikzpicture}[font=\small,>=Stealth,
 box/.style={draw,rounded corners=2pt,align=center,inner sep=5pt,
             minimum height=9mm}]
\node[box,fill=blue!7] (A) at (0,0)
 {Theorem~\ref{thm:main}: flexible sphere};
\node[box] (C) at (0,-1.45)
 {Theorem~\ref{thm:geometric-closure}: geometric closure\quad
  (\S\ref{sec:geometric-closure})};
\node[box,fill=blue!7] (B) at (0,-2.9)
 {Theorem~\ref{thm:annulus-intro} $=$
  Theorem~\ref{thm:new-annular-height}: annular height with flat ends};
\node[box] (N) at (0,-4.35)
 {Nash--Moser branch $q\mapsto\bx(q)$\quad
  (\S\ref{sec:solving-the-equation})\quad and endpoint gain\quad
  (\S\ref{sec:final-proof})};
\node[box] (I) at (0,-5.8)
 {Proposition~\ref{prop:rewrite-volume-inverse}: inverse of the
  corrected linearization near the reference\quad
  (\S\ref{sec:stable-inverse})};
\node[box] (R) at (-4.1,-7.5)
 {Proposition~\ref{prop:rad-fixed-theorem}:\\
  reference profile $w_*$, seed $\phi_*$,\\
  inverse of $\mathscr L_{w_*}$\quad
  (\S\S\ref{sec:rosen-reference-packet}--\ref{sec:radial-reference-inverse})};
\node[box] (F) at (4.1,-7.5)
 {Corrected linearization\\
  on a state-independent domain\\
  (\S\ref{sec:scalar-core}, \S\ref{sec:fixed-geometry})};
\draw[->] (C)--(A);
\draw[->] (B)--(C);
\draw[->] (N)--(B);
\draw[->] (I)--(N);
\draw[->] (R)--(I);
\draw[->] (F)--(I);
\end{tikzpicture}
\caption{The structure of the proof. An arrow means ``is used to
prove''.}
\label{fig:proof-map}
\end{figure}

\begin{table}[htbp]
\centering
\begin{tabularx}{\linewidth}{@{}l Y@{}}
\toprule
Section & What it provides\\
\midrule
\S\ref{sec:geometric-closure} &
 Theorem~\ref{thm:annulus-intro} implies Theorem~\ref{thm:main}:
 planar development, the marked cap, noncongruence, invariants.\\
\S\ref{sec:scalar-core} &
 The Darboux equation as the curvature of a connection form; its
 linearization and conservation law; why the ends must oscillate; the
 two resonances in a finite model; the first residual correction.\\
\S\ref{sec:rosen-reference-packet} &
 Statement of what the reference profile must provide
 (Proposition~\ref{prop:rad-fixed-theorem}); the two solvable models
 from which it is built.\\
\S\ref{sec:radial-tails} &
 The infinite chain of cells at each end; the mass--resistance
 recurrence and the decaying solution; choice of the cell heights.\\
\S\ref{sec:radial-reference-inverse} &
 The inverse of $\mathscr L_{w_*}$ on the whole annulus; proof of
 Proposition~\ref{prop:rad-fixed-theorem}.\\
\S\ref{sec:fixed-geometry} &
 Weighted spaces; the unknowns of the nonlinear problem; the second
 residual correction; coordinates in which every nearby linearization
 has the same transmission conditions.\\
\S\ref{sec:stable-inverse} &
 The inverse of the corrected linearization at nearby states, with tame
 estimates (Proposition~\ref{prop:rewrite-volume-inverse}).\\
\S\ref{sec:solving-the-equation} &
 The coupled linear system and Hamilton's theorem: a smooth branch of
 solutions.\\
\S\ref{sec:final-proof} &
 Improved decay of the metric corrections; the precise annular theorem
 (Theorem~\ref{thm:new-annular-height}); proof of
 Theorem~\ref{thm:main}.\\
Appendix~\ref{app:app-obstructions} &
 A side remark: no finite angular ansatz can solve the nonlinear
 equation.\\
Appendix~\ref{app:dependencies} &
 The external results in the form used, and the order in which the
 constants are chosen.\\
Appendix~\ref{app:notation} & Index of notation.\\
\bottomrule
\end{tabularx}
\caption{Contents by section.}
\label{tab:sections}
\end{table}

\emph{Part~\ref{part:geometry}} (Sections~\ref{sec:geometric-closure}
and~\ref{sec:scalar-core}) contains the geometry and the derivation of
the scalar equation. It is self-contained and elementary.

\emph{Part~\ref{part:reference}}
(Sections~\ref{sec:rosen-reference-packet}--\ref{sec:radial-reference-inverse})
is one-dimensional. After separation of angular modes, the linear
equation \eqref{eq:intro-linear} becomes a family of ordinary
differential equations, and the whole part is devoted to the
construction of one profile $w_*$ for which these equations have the
right kernel and a controlled inverse. Its only output is
Proposition~\ref{prop:rad-fixed-theorem}; a reader willing to accept
that proposition can go directly from Section~\ref{sec:radial-what}
to Part~\ref{part:nonlinear}.

\emph{Part~\ref{part:nonlinear}}
(Sections~\ref{sec:fixed-geometry}--\ref{sec:final-proof}) solves the
nonlinear equation near the reference and proves the two theorems.

The estimates are arranged so that precision is spent where the
decisions above are implemented: the matching of infinitely many
intervals, the common domain of the perturbed operators, and the decay
at the endpoints are proved in detail, while the spectra of the models
and the elementary operator theory are quoted. The specific numbers
that appear in the construction, such as the exponents $2n^3/3$ and
$n^{5/2}$, the power $12$ and the weight exponent $3/4$, are each
introduced together with the inequalities they have to satisfy; none
of them is optimal, and Table~\ref{tab:endpoint-budget} collects them.

\subsection{Context}

\paragraph{Continuous families versus pairs.}
Ghomi's problem \cite[Problem~1.1]{Ghomi2019} asks for a continuous
family. Replacing a dimple by its reflection across a flat rim gives
two different isometric realizations of one metric, but no path between
them. Theorem~\ref{thm:main} provides such a path in the smooth
category.

\paragraph{Infinitesimal versus finite bendings.}
A normal vector field supported inside a planar patch has zero first
variation of the metric. This observation already appears in
Cohn-Vossen's discussion of closed infinitesimally nonrigid surfaces
\cite[footnote~2]{CohnVossen1930}; it does not solve the nonlinear
isometry equations. Efimov's survey \cite{Efimov1948Survey} and his
work on deformations at points of zero curvature
\cite{Efimov1948Zero} belong to the historical setting of this
distinction, and the later survey
\cite[Sections~1.1--1.4]{IvanovaSabitov1991} treats finite bendings,
their infinitesimal equations and higher-order extension. Our starting
point is also an infinitesimal bending, the solution of
\eqref{eq:intro-linear}; the analysis of
Parts~\ref{part:reference} and~\ref{part:nonlinear} is what turns it
into an actual height.

\paragraph{Rigidity.}
On the rigid side, hypotheses on curvature and on boundary data are
decisive. For example, Ghomi and Spruck
\cite[Theorem~1.1]{GhomiSpruck2020} prove that two properly oriented
$C^{2,1}$ positively curved isometric immersions which agree on a curve
segment agree everywhere on their connected domain. The annulus of
Theorem~\ref{thm:annulus-intro} has curvature of both signs, with
infinitely many transitions at each flat end, so neither strict
positivity nor a finite collection of nondegenerate interfaces
describes its geometry.

\paragraph{Other categories.}
Connelly's flexible embedded polyhedral sphere \cite{Connelly1977}
establishes flexibility for polyhedra. Rounding its edges would change
the metric, so it does not furnish a smooth bending. The $C^1$ theorems
of Nash and Kuiper \cite{Nash1954,Kuiper1955} concern another
regularity regime; in the smooth problem the second fundamental form
and the type of the Darboux equation must be controlled at every scale.

\paragraph{Smooth isometric embedding.}
Our analytic method is closest to the theory of smooth isometric
embedding. Smoothing already compensates for the loss of derivatives in
Nash's argument \cite{Nash1956}. Han, Hong and Lin
\cite[Theorem~1.1]{HanHongLin2003} use a Nash--Moser--H\"ormander
scheme for nonpositive curvature under a specified degeneracy
condition, and Han and Khuri treat local problems of mixed type with
finite-order changes of type \cite{HanKhuri2010,HanKhuri2013}. In
particular, the modification of the linearization by a multiple of the
residual in \cite[pp.~650--653]{HanKhuri2010} is a precedent for our
residual corrections. The new issue here is the simultaneous control of
infinitely many interfaces and of both flat ends on a single closed
domain.

\paragraph{The invariants.}
The three invariants in Theorem~\ref{thm:main} follow directly from the
form \eqref{eq:phase-family} of the motion. They should be compared
with general results of different scope: Sabitov's volume theorem
\cite{Sabitov1996} is polyhedral, and Almgren and Rivin
\cite{AlmgrenRivin1998} prove the invariance of the mean-curvature
integral under smooth bendings. Motions in which curvature precesses
around a shell also arise in the mechanical model of Hamouche, Maurini,
Vidoli and Vincenti \cite{Hamouche2017}; that comparison supplies
neither the flat collars nor the closed embedded surface.

\subsection{Conventions}
\label{subsec:conventions}

Throughout, $\T=\R/2\pi\Z$, the angular average of a function is
\[
 \av f=\frac1{2\pi}\int_0^{2\pi}f\,d\theta ,
\]
and $R_\vartheta$ denotes a rotation by the angle $\vartheta$, of the
plane, of $\R^3$ about the vertical axis, or of the circle
($R_\vartheta\theta=\theta+\vartheta$), as the context requires. A
prime denotes a derivative with respect to the radius $r$ unless
another variable is indicated by a subscript.

The letters $C$ and $c$ denote positive constants, large and small
respectively, whose values may change from line to line; a subscript
indicates dependence, as in $C_s$. We write $a\asymp b$ when
$cb\le a\le Cb$. Constants that are fixed once and referred to later
carry a distinguishing mark, as in $C_*$, $c_0$ or $N$.

The paper uses a fairly large number of symbols. Each of
Parts~\ref{part:reference} and~\ref{part:nonlinear} begins by fixing
its vocabulary, and Appendix~\ref{app:notation} is an index of
notation. Some letters have a local meaning inside a single proof;
this is always said where it happens.
